% ---------------------------------------------------------------------------
% Macros du manuscrit. Les chapitres n'en définissent AUCUNE autre.
% ---------------------------------------------------------------------------

% Racine du dépôt, relative à docs/manuscrit (pour \rustsnippet & co).
\newcommand{\repoRoot}{../..}
% Commit photographié par le manuscrit.
\newcommand{\snapshotCommit}{e67b10a}
\newcommand{\snapshotDate}{27 septembre 2026}

% ---- Texte -----------------------------------------------------------------
% Identifiant / chemin en machine à écrire, accepte _ # & % $ ^ ~ sans
% échappement ET les caractères UTF-8 (⊤, ✓, é…). Principe : on détokenise
% l'argument, puis on le relit (\scantokens) avec les caractères spéciaux
% rendus inertes, tandis que les octets UTF-8 restent actifs et sont décodés
% par inputenc. \{ \} \# \ldots restent des commandes et s'impriment.
\makeatletter
\newcommand{\ra@codecatcodes}{%
  \catcode`\_=12 \catcode`\#=12 \catcode`\&=12 \catcode`\%=12
  \catcode`\$=12 \catcode`\^=12 \catcode`\~=12 \endlinechar=-1 }
\DeclareRobustCommand{\code}[1]{\begingroup\ra@codecatcodes\edef\ra@tmp{\detokenize{#1}}\texttt{\scantokens\expandafter{\ra@tmp}}\endgroup}
\DeclareRobustCommand{\file}[1]{\begingroup\ra@codecatcodes\edef\ra@tmp{\detokenize{#1}}\texttt{\scantokens\expandafter{\ra@tmp}}\endgroup}
% Dans les signets PDF (hyperref), \code et \file rendent leur argument tel quel.
\pdfstringdefDisableCommands{\let\code\@firstofone \let\file\@firstofone}
\makeatother
% Référence à un endroit du code : \src{src/ir/expr.rs}{12}{40}  /  \srcline{src/ir/expr.rs}{12}
\newcommand{\src}[3]{\file{#1}, l.~#2--#3}
\newcommand{\srcline}[2]{\file{#1}, l.~#2}
% Nom d'une règle du catalogue : \regle{infinite-loop}
\newcommand{\regle}[1]{\textsf{\detokenize{#1}}}
% Nom d'une relation du moteur : \relation{writers}
\newcommand{\relation}[1]{\textsf{\detokenize{#1}}}
% Référence à une ADR : \adr{27} -> lien vers l'annexe, si le label existe
\newcommand{\adr}[1]{\hyperref[adr:#1]{ADR-#1}}
% Issue GitHub : \issue{63}
\newcommand{\issue}[1]{\href{https://github.com/rboudrouss/reactant-analyzer/issues/#1}{\##1}}
% Terme défini pour la première fois (italique + index)
\newcommand{\terme}[1]{\emph{#1}\index{#1}}
% Entrée d'index sans mise en forme
\newcommand{\idx}[1]{\index{#1}}
% Niveaux de diagnostic
\newcommand{\Error}{\textsc{Error}}
\newcommand{\Warning}{\textsc{Warning}}
\newcommand{\Info}{\textsc{Info}}
% Nom du projet
\newcommand{\reactant}{\textsf{reactant-analyzer}}
\newcommand{\React}{React}

% ---- Mathématiques ---------------------------------------------------------
\newcommand{\must}{\mathsf{must}}
\newcommand{\may}{\mathsf{may}}
\newcommand{\Top}{\top}
\newcommand{\Bot}{\bot}
\newcommand{\join}{\sqcup}
\newcommand{\meet}{\sqcap}
\newcommand{\widen}{\mathbin{\nabla}}
\newcommand{\lleq}{\sqsubseteq}
\newcommand{\lgeq}{\sqsupseteq}
\newcommand{\lfp}{\mathrm{lfp}}
\newcommand{\gam}{\gamma}
\newcommand{\abs}{\alpha}
\newcommand{\sem}[1]{\llbracket #1 \rrbracket}
\providecommand{\llbracket}{[\![}
\providecommand{\rrbracket}{]\!]}
\newcommand{\Int}{\mathbb{Z}}
\newcommand{\Nat}{\mathbb{N}}
\newcommand{\Bool}{\mathbb{B}}
\newcommand{\powerset}{\mathcal{P}}
\newcommand{\Stab}{\mathcal{S}}
\newcommand{\Val}{\mathcal{V}}
\newcommand{\Env}{\mathcal{E}}
\newcommand{\Heap}{\mathcal{H}}
\newcommand{\Store}{\mathcal{M}}
\newcommand{\Render}{\mathsf{render}}
\newcommand{\Commit}{\mathsf{commit}}
\newcommand{\Effect}{\mathsf{effect}}
\newcommand{\Tick}{\mathsf{tick}}
\newcommand{\Fix}{\mathsf{fix}}
\DeclareMathOperator{\dom}{dom}
\DeclareMathOperator{\reach}{reach}
\DeclareMathOperator{\pred}{pred}
\DeclareMathOperator{\succs}{succ}

% ---- Extraits de code pris dans le dépôt -----------------------------------
% \rustsnippet[options listings]{chemin relatif au dépôt}{première ligne}{dernière ligne}{légende}
% Inclut VERBATIM les lignes du fichier, numérotées avec leurs vrais numéros.
\newcommand{\rustsnippet}[5][]{%
  \lstinputlisting[language=Rust,firstline=#3,lastline=#4,firstnumber=#3,
    caption={\file{#2}, l.~#3--#4 --- #5},#1]{\repoRoot/#2}}
% Même chose pour un fichier TSX/TS/JS du dépôt (fixtures, tests)
\newcommand{\tsxsnippet}[5][]{%
  \lstinputlisting[language=TypeScript,firstline=#3,lastline=#4,firstnumber=#3,
    caption={\file{#2}, l.~#3--#4 --- #5},#1]{\repoRoot/#2}}
% Fichier entier
\newcommand{\rustfile}[3][]{\lstinputlisting[language=Rust,caption={\file{#2} --- #3},#1]{\repoRoot/#2}}
\newcommand{\tsxfile}[3][]{\lstinputlisting[language=TypeScript,caption={\file{#2} --- #3},#1]{\repoRoot/#2}}
% Sortie d'une commande (texte brut, sans numéros) : \begin{sortie} ... \end{sortie}
\lstnewenvironment{sortie}[1][]{\lstset{language={},numbers=none,basicstyle=\ttfamily\scriptsize,backgroundcolor=\color{white},frame=leftline,#1}}{}
% Terminal
\lstnewenvironment{shell}[1][]{\lstset{language=bash,numbers=none,basicstyle=\ttfamily\footnotesize,#1}}{}

% ---- Structure de chapitre -------------------------------------------------
% Encadré d'objectifs en tête de chapitre
\newenvironment{objectifs}{%
  \par\noindent\begin{minipage}{\linewidth}\small
  \textbf{Objectifs du chapitre.}\begin{itemize}}{\end{itemize}\end{minipage}\par\bigskip}
% Résumé de fin de chapitre
\newenvironment{resume}{\section*{Résumé du chapitre}\addcontentsline{toc}{section}{Résumé du chapitre}}{}
% Prérequis
\newcommand{\prerequis}[1]{\par\noindent{\small\textbf{Prérequis :} #1}\par\medskip}

% ---- Figures ---------------------------------------------------------------
% Style TikZ commun (si TikZ est chargé)
\ifhastikz
\tikzset{
  bloc/.style={draw=ra-blue,fill=ra-bg,rounded corners=2pt,thick,align=center,inner sep=5pt,font=\small},
  noeud/.style={draw=ra-blue,circle,thick,inner sep=2pt,font=\small},
  cfgbloc/.style={draw=ra-gray,fill=white,rectangle,align=left,inner sep=4pt,font=\ttfamily\scriptsize},
  fleche/.style={-{Stealth[length=2mm]},thick,draw=ra-gray},
  flechemust/.style={fleche,draw=ra-blue},
  flechemay/.style={fleche,dashed,draw=ra-amber},
  treillis/.style={font=\small},
  etiquette/.style={font=\scriptsize\itshape,text=ra-gray},
}
\fi
